% ============================================================
% Appendix Z.5 — FRFP Case Study (Real Estate / Finance):
% Zillow Offers and Algorithmic Home-Flipping
% ============================================================

\section{FRFP Case Study (Real Estate / Finance): Zillow Offers}
\label{app:frfp-zillow-offers}

This appendix analyses the rise and collapse of Zillow’s iBuying business,
\emph{Zillow Offers}, as a case study for FRFP in a commercial, quasi-financial
setting.
Rather than treating it as a one-off business mishap, we examine how its
structure interacted with FRFP’s explicit/tacit split, boundary design, safe
horizons, and institutional semantics.

\subsection{Background}

Zillow is best known as an online real estate marketplace and for its
``Zestimate'' home-value estimates.
In 2018 it launched Zillow Offers, an ``iBuying'' programme in which Zillow
itself would:

\begin{itemize}
  \item make instant cash offers to homeowners using algorithmic valuations;
  \item purchase homes directly, perform light renovations, and resell them;
  \item profit from the spread between purchase price, renovation costs,
        and resale price.
\end{itemize}

Zillow framed its machine-learning models as a competitive advantage:
they claimed to be able to price homes with high accuracy and to do so at
scale, turning their Zestimate and related models into an internal
trading engine.

By 2021, Zillow Offers was purchasing thousands of homes per quarter.
In Q3 2021 alone, Zillow bought roughly 9,700 homes, many at prices that
subsequently proved to be too high.
When market conditions shifted and Zillow attempted to resell, the company
faced large losses.
In November 2021, Zillow announced it would shut down Zillow Offers, record
more than \$500M in write-downs, and lay off about 25\% of its workforce.

Analyses highlight several proximate causes:

\begin{itemize}
  \item valuation models that systematically overpaid in certain markets,
        especially during volatile conditions;
  \item feedback and selection effects: homeowners with overvalued or
        problematic properties were more likely to accept Zillow’s offers;
  \item operational and scaling challenges in renovating and reselling homes
        at volume;
  \item leadership and governance decisions that concentrated risk
        based on confidence in the models.
\end{itemize}

In FRFP terms, Zillow Offers is a clean example of AI-driven decision-making
where explicit models and pipelines are tightly coupled to financial and
operational commitments, with human and institutional actors shaping and
interpreting them.

\subsection{Explicit and Tacit Spaces in Zillow Offers}

\subsubsection*{Explicit space \(E_{\mathrm{zillow}}\)}

The explicit domain for Zillow Offers includes:

\begin{itemize}
  \item structured property features:
        location, square footage, bed/bath counts, photos, time-on-market,
        local comparable sales, etc.;
  \item the valuation models:
        updated Zestimate architectures (eventually based on neural networks)
        and internal purchase-price models tuned for flipping;
  \item decision rules that map model outputs to offers:
        e.g.\ target discount or premium vs.\ predicted resale price, margin
        buffers for renovation and transaction costs;
  \item operational artefacts:
        purchase offers, contracts, renovation plans, and listing prices;
  \item dashboards and risk metrics tracking inventory, margins,
        and exposure by market.
\end{itemize}

All of these are explicit artefacts in \(E_{\mathrm{zillow}}\): they can be
stored, versioned, and manipulated by automated systems.

\subsubsection*{Tacit space \(T_{\mathrm{zillow}}\)}

Tacit space comprises the epistemic states of:

\begin{itemize}
  \item data scientists and engineers:
        their informal sense of where models work or fail; their understanding
        of model assumptions and the limits of generalisation;
  \item local real estate experts:
        tacit knowledge of neighbourhood idiosyncrasies, construction quality,
        market sentiment, and micro-location effects;
  \item executives and product leaders:
        their beliefs about strategic risk appetite, market dynamics, and how
        much to trust the models;
  \item investors and board members:
        their high-level view of risk, growth, and business viability.
\end{itemize}

These tacit states guide how explicit model outputs are interpreted, when
they are trusted, and how aggressively they are used to commit capital.

\subsection{Boundary Design: Algorithmic Prices vs.\ Human Oversight}

\subsubsection*{Intended boundary}

Conceptually, Zillow positioned its models as sophisticated tools in a
human-led business:

\begin{itemize}
  \item algorithms generate valuations and offers;
  \item humans set strategy and oversee operations;
  \item market discipline (ability to resell at a profit) validates the
        whole pipeline.
\end{itemize}

In FRFP terms, the intended boundary is:

\begin{itemize}
  \item AI systems in \(E_{\mathrm{zillow}}\) produce valuations and
        suggested offers;
  \item human managers in \(T_{\mathrm{zillow}}\) endorse or adjust these
        offers before they become binding commitments.
\end{itemize}

\subsubsection*{De facto boundary: models as capital allocation engines}

In practice, several reports suggest that Zillow Offers operated closer to a
semi-automated capital allocation engine:

\begin{itemize}
  \item offers were often based on model outputs plus relatively simple
        rules, with limited deal-by-deal human adjustment;
  \item scale pressures meant many offers were generated and accepted with
        minimal local expert intervention, especially in hot markets;
  \item internal confidence in the models (and public messaging about
        algorithmic advantage) encouraged leadership to ramp up volume and
        risk exposure.
\end{itemize}

Thus, the effective boundary shifted toward:

\begin{itemize}
  \item model-driven offer generation and acceptance as de facto capital
        commitments, subject to only thin human review;
  \item human oversight focused more on macro metrics and scaling decisions
        than on granular corrections of model outputs.
\end{itemize}

From an FRFP perspective, this is a form of \emph{boundary capture}:
correctness about price and risk migrates from tacit human judgment into
explicit pipelines.

\subsection{Protocol-Level Failures in FRFP Terms}

\subsubsection*{Z.5.1 Explicit-Space Model Limitations}

\paragraph{Problem.}
Multiple analyses, including academic case studies and industry writeups,
highlight limitations in Zillow’s explicit valuation models:

\begin{itemize}
  \item Models were trained to estimate market values based on historical
        sales, but the iBuying use case required predicting \emph{future}
        sale prices in volatile markets for specific holding periods.
  \item The algorithm appears to have systematically overpaid in certain
        rapidly rising or changing markets, leading to negative expected
        margins when conditions shifted.
  \item Public commentary suggests selection effects:
        homeowners with hidden defects or less desirable homes were more
        likely to accept Zillow’s offers, while owners of highly desirable
        homes had better alternatives.
\end{itemize}

In FRFP language, the explicit mapping
\[
  f_{\mathrm{offer}} : \text{property features, market state} \to
  \text{offer price}
\]
embeds structural mismatches:

\begin{itemize}
  \item it does not fully capture tacit factors known to local agents
        (e.g.\ construction quality, micro-location quirks, noisy seller
        motives);
  \item it assumes stability in relations between features and prices that
        fails under regime shifts and high volatility.
\end{itemize}

\paragraph{Consequence.}
Even with perfect boundary discipline, reliance on such a model at large
scale would produce mispriced portfolios.
The FRFP kernel identifies this as a defective explicit semantics feeding
into the decision boundary.

\subsubsection*{Z.5.2 Boundary Design: From Tool to Trader}

\paragraph{Problem.}
Zillow Offers effectively elevated its models from \emph{advisory tools}
to \emph{traders}:

\begin{itemize}
  \item offers went out at scale, often automatically, with human review
        focused on exceptions or special cases;
  \item once an offer was accepted by a seller, Zillow incurred substantial
        capital commitments and operational obligations;
  \item governance practices appear to have focused on aggregate exposure,
        not on continuous FRFP-style evaluation of boundary behaviour.
\end{itemize}

In FRFP terms:

\begin{itemize}
  \item the boundary between explicit offers and tacit acceptance was
        located at the algorithmic offer stage, not at a human review stage;
  \item human endorsement happened upstream, at the level of approving the
        \emph{use of the model} and the operating policy, rather than on
        individual decisions;
  \item this boundary placement effectively automated much of the correctness
        function for pricing.
\end{itemize}

\paragraph{Consequence.}
When models were wrong, they were wrong systematically and at scale.
Safe horizons were defined more by capital constraints and external shocks
than by deliberate FRFP-style limits on how many model-driven acquisitions
could proceed before re-grounding.

\subsubsection*{Z.5.3 Absent Safe Horizons and Re-Grounding}

\paragraph{Problem.}
Zillow Offers operated as an ongoing, high-volume programme:

\begin{itemize}
  \item thousands of homes purchased and sold across markets;
  \item models and policies updated, but within a continuous, expanding
        pipeline;
  \item no clear public evidence of hard limits on how many homes could be
        acquired under a given model configuration before thorough review.
\end{itemize}

FRFP’s dynamics suggest that:

\begin{itemize}
  \item long-horizon reliance on explicit models without structured
        re-grounding invites degradation of tacit risk awareness;
  \item decision-makers may come to over-trust dashboards and backtests;
  \item model failure under new conditions (e.g.\ sudden housing-market
        inflection) is almost inevitable.
\end{itemize}

\paragraph{Consequence.}
When the U.S.\ housing market began to shift in 2021, Zillow’s model-driven
acquisitions left it holding inventory bought at prices that proved too high,
leading to more than \$500M in write-downs and the shuttering of the unit.
From an FRFP viewpoint, this is a classic safe-horizon violation:
the depth of model-driven commitments exceeded the robustness of explicit
semantics and tacit oversight.

\subsubsection*{Z.5.4 Institutional Operator Choices}

\paragraph{Problem.}
Leadership decisions played a central role:

\begin{itemize}
  \item CEO Rich Barton publicly stated that the failure was due, in part, to
        the AI’s inability to predict prices accurately enough.
  \item Other analyses emphasise that leadership chose to scale rapidly,
        accept thin margins, and concentrate risk based on confidence in the
        modelling approach, rather than on cautious, incremental pilots.
\end{itemize}

In FRFP terms, the institutional operator \(G_{\mathrm{zillow}}\)---the
mapping from tacit states of executives, risk managers, and boards to new
tacit states and policies---aligned strongly with algorithmic optimism:

\begin{itemize}
  \item explicit model performance metrics were treated as decisive evidence
        of readiness;
  \item tacit local expertise (real estate agents, local partners) appears to
        have had limited structural authority in gating model-driven offers;
  \item there was no robust institutional mechanism to throttle or reverse
        scaling in response to early signs of mispricing until large losses
        materialised.
\end{itemize}

\paragraph{Consequence.}
FRFP would classify this as an instance of explicit-only governance:
institutional decisions are driven by metrics and dashboards produced by
explicit pipelines, without a sufficiently strong, structurally distinct
tacit governance layer capable of enforcing conservative boundaries and
finite safe horizons.

\subsection{Counterfactual: An FRFP-Aligned iBuying Architecture}

What might an FRFP-conforming iBuying programme have looked like?

\subsubsection*{Z.5* Explicit/Tacit Role Separation}

\begin{itemize}
  \item Models are designated as \emph{valuation assistants}:
        they propose price bands and risk indicators, but do not generate
        binding offers autonomously.
  \item Local human experts and central risk committees retain explicit
        authority to set offers and decide volume limits per market.
\end{itemize}

\subsubsection*{Z.5* Boundary Design for Offers}

\begin{itemize}
  \item Each offer above a certain dollar or risk threshold is treated as a
        boundary artefact:
        \begin{itemize}
          \item it must be explicitly endorsed by a human operator (local
                specialist or central desk) who sees both model outputs and
                key raw features (photos, inspections, comps);
          \item the endorsement and rationale are logged in a decision ledger.
        \end{itemize}
  \item Automated offers may be permitted only within a tightly constrained
        low-risk band (e.g.\ low-priced homes with strong liquidity and
        conservative discounts).
\end{itemize}

\subsubsection*{Z.5* Safe Horizons and Pilot Design}

\begin{itemize}
  \item The programme is structured around finite-horizon pilots:
        \begin{itemize}
          \item a limited number of homes per market and per model version;
          \item explicit triggers (e.g.\ loss thresholds, variance thresholds)
                that force a pause and review;
          \item mandatory re-grounding in tacit local expertise after each
                pilot episode.
        \end{itemize}
  \item Safe horizons are shorter in volatile markets and extended only after
        successful, audited performance.
\end{itemize}

\subsubsection*{Z.5* Institutional Governance}

\begin{itemize}
  \item A cross-functional governance body (data science, risk, operations,
        local real estate, finance) operates as an FRFP institutional operator:
        \begin{itemize}
          \item it defines admissible model uses and boundary rules;
          \item it sets and adjusts safe horizons by market and product;
          \item it monitors performance, selection effects, and
                model-environment mismatch.
        \end{itemize}
  \item Major scaling decisions (e.g.\ moving from pilot to thousands of
        homes per quarter) require explicit sign-off with documented
        reasoning about model limitations and macro risk.
\end{itemize}

Under such an architecture, Zillow might still have faced model error and
market risk, but those errors would have been bounded by explicit FRFP
constraints:
fewer homes mispriced at a time, more opportunities for human re-grounding,
and clearer responsibility for boundary decisions.

\subsection{Lessons for FRFP and Commercial AI}

The Zillow Offers collapse reinforces several FRFP themes in a commercial,
non-regulated domain:

\begin{enumerate}
  \item \textbf{Predictive accuracy is not the same as decision correctness.}
        Models can be reasonably accurate on historical data yet still fail
        when used to drive large-scale, capital-intensive decisions in
        changing environments.

  \item \textbf{Boundary placement is a strategic risk decision.}
        Treating model outputs as final offers effectively moves the boundary
        into explicit space.
        FRFP recommends keeping correctness at human boundaries with clear,
        logged endorsements.

  \item \textbf{Safe horizons are crucial for AI-driven trading and pricing.}
        Committing to thousands of model-driven purchases without hard caps,
        episodic pilots, and re-grounding is structurally risky.

  \item \textbf{Institutional operators shape technical risk.}
        Governance choices about scaling, margin targets, and reliance on
        models are as important as model architecture.
        FRFP provides a way to describe those choices as operations on tacit
        states, subject to non-automation constraints.

  \item \textbf{Formal architectures clarify postmortems.}
        Rather than attributing Zillow’s failure solely to ``bad AI'' or
        ``unlucky market timing'', FRFP highlights a pattern:
        imperfect explicit models, boundary capture, absent safe horizons,
        and explicit-only governance combined to create a structurally
        fragile system.
\end{enumerate}

In this light, Zillow Offers is not just an AI business gone wrong, but a
textbook example of why FRFP-style architectures and governance are needed
whenever explicit AI models are coupled tightly to large-scale financial and
operational commitments.
